\subsection{Domain Adaptation with CORAL}

After constructing a shared label space, we apply CORAL~\cite{sun2016coral}
to reduce distribution divergence between PlantVillage and PlantDoc feature
representations. CORAL aligns the second-order statistics of source and target
feature distributions by minimizing the squared Frobenius norm of the
difference between their covariance matrices~\cite{sun2015return}, computed
from activations of the penultimate layer of a ResNet-50 backbone.

Let $\mathbf{C}_S$ and $\mathbf{C}_T$ denote the covariance matrices of the
source and target features, each of dimension $d \times d$ (with $d = 2048$
for ResNet-50). The CORAL alignment loss is:
%
\begin{equation}
  \mathcal{L}_{\mathrm{CORAL}}
    = \frac{1}{4d^{2}}
      \left\lVert \mathbf{C}_S - \mathbf{C}_T \right\rVert_F^{2}\,,
  \label{eq:coral}
\end{equation}
%
where $\lVert \cdot \rVert_F$ is the Frobenius norm. The overall training
objective combines supervised cross-entropy on the source data with the
CORAL term:
%
\begin{equation}
  \mathcal{L}
    = \mathcal{L}_{\mathrm{CE}}
    + \lambda \, \mathcal{L}_{\mathrm{CORAL}}\,,
  \label{eq:total_loss}
\end{equation}
%
where $\lambda$ controls the strength of distribution regularization relative
to classification accuracy.

The choice of $\lambda$ reflects a bias--variance trade-off. Small values make
the covariance-alignment signal too weak to overcome the structural mismatch
between studio and field images; the adapted model behaves almost identically
to the non-adapted one. Very large values force the model to match feature
covariances so aggressively that discriminative but source-specific information
is suppressed---for example, color-saturation cues for fungal lesions that
appear differently under natural versus artificial lighting. At
$\lambda = 10.0$, we observe a consistent slight degradation in target~F1
across all three labeling strategies compared with $\lambda = 1.0$, consistent
with this over-regularization mechanism.

We perform CORAL adaptation using a ResNet-50 backbone pretrained on ImageNet.
The classification head is retrained from scratch to match the size of the
shared class set specific to each labeling strategy. Source features are taken
from the penultimate global average pooling layer (dimension $d = 2048$). The
batch size is 32 in all cases; target-domain images are included in each batch
to estimate~$\mathbf{C}_T$, but they do not contribute to the classification
loss~$\mathcal{L}_{\mathrm{CE}}$. All CORAL runs use
AdamW~\cite{loshchilov2019adamw} with a learning rate of $3 \times 10^{-4}$,
weight decay $10^{-4}$, and are trained for 20~epochs without a learning-rate
scheduler.
